\documentclass[11pt,a4paper]{article}
\usepackage[margin=0.8in]{geometry}
\usepackage[T1]{fontenc}
\usepackage[utf8]{inputenc}
\usepackage{hyperref}
\usepackage{booktabs}
\usepackage{longtable}
\usepackage{array}
\usepackage{enumitem}
\usepackage{xcolor}
\usepackage{verbatim}
\hypersetup{colorlinks=true, linkcolor=blue!50!black, urlcolor=blue!50!black}
\setlist[itemize]{leftmargin=*, itemsep=2pt, topsep=2pt}
\setlist[enumerate]{leftmargin=*, itemsep=2pt, topsep=2pt}
\newcommand{\code}[1]{\texttt{#1}}
\newcommand{\repo}[1]{\texttt{#1}}
\newcolumntype{P}[1]{>{\raggedright\arraybackslash}p{#1}}

\title{Business Study and Assumption Validation Report\\\large Digital Procurement and PLS Seedbed}
\author{Senior Software Architecture, Business Analysis, and Technical Research Review}
\date{2026-05-27}

\begin{document}
\maketitle

\begin{center}
\begin{tabular}{ll}
Repository & \repo{raichiiiiiii/Blockchain-Based-E-Procurement-System} \\
Branch inspected & \repo{origin/main} \\
Baseline commit observed & \repo{679c2e18b5894d9df3bdbea9a20d403610bb69b4} \\
Report branch & \repo{docs/business-study-assumption-validation} \\
Report status & First validation draft for review before merge \\
\end{tabular}
\end{center}

\section*{Executive Summary}

The current \repo{origin/main} is aligned with the documented direction of a compliance-first digital procurement evidence platform with blockchain proof anchoring and Shariah-governed PLS seedbed support. It is not aligned with a production-certified or commercial-ready claim. Repository evidence states the current maturity as supervisor-demo plus selected pilot-hardening features, not commercial-ready or production-certified. The latest production-extension validation records completed hardening across payment adapters, escrow release/dispute, IoT/QR/EPCIS intake, document handling, Shariah certificate artifacts, signed export bundles, ERP/accounting adapter foundation, contract negotiation, external API gateway, containerized deployment, observability, and backup/restore runbooks.

The largest remaining gap is not feature count. It is validation depth. Several implemented capabilities are intentionally adapter foundations, local simulations, or metadata-only controls. The system still needs a formal assumption register, standard-to-domain mapping closure, persistence completion beyond the core Postgres-backed subset, production-like Fabric lab validation, managed secret/key operations, legal/Shariah stakeholder review, and pilot-readiness evidence before development continues toward a fully deployable system.

\section{Repository Inspection Summary}

\subsection{What exists}

The inspected repository contains:

\begin{itemize}
  \item Source code under \repo{src/app} and \repo{src/modules} using TypeScript, Fastify, React/Vite frontend, module-level application/domain/infrastructure boundaries, and Postgres adapters for selected runtime repositories.
  \item Deployment files including \repo{Dockerfile.backend}, \repo{Dockerfile.frontend}, \repo{docker-compose.yml}, and \repo{docker-compose.app.yml}.
  \item Architecture and contract documentation under \repo{docs/architecture}, \repo{docs/contracts}, \repo{docs/process}, and \repo{docs/runbooks}.
  \item QA/evidence records under \repo{docs/evidence/qa}, including final release candidate and production-extension validation evidence.
  \item Production-extension backlog rows under \repo{backlog/production-extension-roadmap.csv}.
  \item Fabric local chaincode baseline and production-consortium templates under \repo{chaincode} and \repo{fabric/production-consortium}.
\end{itemize}

\subsection{System boundaries observed}

\begin{longtable}{P{0.20\linewidth}P{0.36\linewidth}P{0.32\linewidth}}
\toprule
Boundary & Observed responsibility & Validation comment \\
\midrule
Application API & Fastify server registers auth, membership, access control, procurement, delivery evidence, KYC/AML, blockchain proof, escrow, reporting/export, PLS, security, operations, external API, ERP/accounting, documents, contracts, and payments. & Broad backend composition exists. Some modules remain in-memory unless explicitly backed by Postgres. \\
PostgreSQL & Operational persistence for auth, membership, roles, access audit, procurement lifecycle, procurement orders, delivery evidence, blockchain metadata, and escrow. & Not yet a complete persistence layer for every production-extension module. \\
Fabric / blockchain & Proof anchoring through AuditAnchor and production-consortium templates. & Production consortium implementation remains open; no real CA/MSP/channel deployment evidence. \\
Frontend & React/Vite product shell, role dashboards, proof/escrow/export/PLS/status surfaces. & Product UI appears aligned with product vocabulary rule; framework replacement is not currently justified. \\
External integrations & External API gateway, ERP/accounting adapter, payment adapter, document/signature, EPCIS/IoT/QR interfaces. & Mostly adapter foundations and local/sandbox implementations; production certification is not evidenced. \\
\bottomrule
\end{longtable}

\subsection{What appears complete}

Completed production-extension rows include PBI-437, PBI-439 to PBI-451, and PBI-453 to PBI-461. The open rows are intentionally limited to PBI-436, PBI-438, PBI-452, and PBI-462. The production-extension validation records passing build, frontend build, tests, database dry-runs, AuditAnchor chaincode checks, Docker compose checks, deployable smoke, UI label scan, and browser smoke across all nine demo accounts.

\subsection{What appears incomplete}

\begin{itemize}
  \item PBI-438 remains incomplete because no live production-like Fabric consortium exists with real CA/MSP material, channel artifacts, lifecycle commit, and cross-organization smoke evidence.
  \item PBI-452 remains open because blockchain status visualization does not yet cover full Fabric network health and all PLS-wide/operator states.
  \item PBI-462 remains open because procurement and public-contract standards mapping research is not closed.
  \item Runtime Postgres coverage is incomplete for the newest extension modules. Several repositories still default to in-memory adapters after the Postgres-backed core repositories are configured.
  \item Legal, Shariah, banking, ERP, logistics, and regulator assumptions remain project assumptions until validated with named stakeholders or controlled external pilots.
\end{itemize}

\subsection{Undocumented or weakly justified areas}

\begin{itemize}
  \item The distinction between \repo{docs/report} and the requested \repo{docs/reports} path should be normalized in a documentation index update.
  \item Some older durable docs still contain Sprint 1 or Aider-era wording and should be refreshed to Codex/current-agent-neutral operating language.
  \item Production operating policies for secrets, keys, retention, legal hold, incident response, and environment promotion are still runbook-level rather than enforced system controls.
\end{itemize}

\section{Documentation Review}

\begin{longtable}{P{0.30\linewidth}P{0.30\linewidth}P{0.28\linewidth}}
\toprule
Document & Purpose & Review finding \\
\midrule
\repo{docs/evidence/qa/PRODUCTION_EXTENSION_RELEASE_VALIDATION.md} & Release evidence for production-extension consolidation. & Primary evidence for current readiness; explicitly limits claims to supervisor demo/internal technical review. \\
\repo{backlog/production-extension-roadmap.csv} & PBI-436 to PBI-462 execution roadmap. & Current extension status source. PBI-462 still planned. \\
\repo{docs/architecture/EXTENDED_PRODUCTION_ARCHITECTURE_PLAN.tex} & Production-extension target architecture and business/technical transition plan. & Aligns with current ports/adapters direction. \\
\repo{docs/architecture/PRODUCTION_FABRIC_CONSORTIUM_ARCHITECTURE.md} & Fabric consortium plan and foundation boundary. & Good architecture plan; explicitly not live consortium implementation. \\
\repo{docs/contracts/ESCROW_WORKFLOW_CONTRACT.md} & Escrow state and transition contract. & Updated with release/dispute pilot-hardening states and no-payment-execution rule. \\
\repo{docs/contracts/BLOCKCHAIN_ANCHOR_CONTRACT.md} & Anchor and proof API semantics. & Good claim boundary: proof anchoring does not replace workflow, authorization, or Postgres state. \\
\repo{docs/process/CODING_RULES.md} & Engineering constraints. & Useful technical rules, but agent-process wording is stale and should be updated. \\
\bottomrule
\end{longtable}

\section{Assumption Register}

\begin{longtable}{P{0.08\linewidth}P{0.32\linewidth}P{0.20\linewidth}P{0.28\linewidth}}
\toprule
ID & Assumption & Status & Evidence / next action \\
\midrule
A01 & Product is supervisor-demo plus selected pilot hardening, not production-certified. & Validated by project documentation. & Production-extension validation and scorecard state this directly. Preserve claim boundary. \\
A02 & Credential-only login and backend-seeded demo accounts are the normal path. & Validated by repository evidence. & Demo accounts and no role-card login are recorded in release validation. \\
A03 & PostgreSQL is the operational system of record. & Partially validated. & Core repositories have Postgres adapters; newer extension modules still need persistent adapters/migrations. \\
A04 & Fabric stores proof metadata and hashes, not raw business data. & Validated by documentation and external Fabric research. & Fabric private data supports private data plus ledger hashes; repository also prohibits raw KYC/documents/payment data on-chain. \\
A05 & Production Fabric consortium is planned but not implemented. & Validated by evidence. & PBI-438 remains open until real CA/MSP/channel/lifecycle/smoke evidence exists. \\
A06 & External standards are adapter/input-output formats, not the internal domain model. & Validated by architecture direction; PBI-462 unresolved. & Must close with standards mapping research before deeper ERP/public-contract implementation. \\
A07 & Payment support is sandbox/manual and does not move real money. & Validated by docs and evidence. & Payment adapters should remain behind PaymentPort until certification and rails are selected. \\
A08 & ISO 20022 support is mapping-only, not bank certification. & Validated by docs; external research supports repository/catalogue discipline. & Do not claim bank execution. \\
A09 & Shariah certification artifacts are record tracking, not formal external certification. & Validated by evidence. & Requires named Shariah/legal review before pilot claim. \\
A10 & Document signature verification is local/adapted metadata, not legal e-signature certification. & Validated by project evidence. & Requires external trust store/legal signature policy later. \\
A11 & Containerized deployment is internal-pilot oriented, not hardened production hosting. & Validated by deployment evidence. & Needs secrets, TLS, image hardening, environment promotion, and backup drills. \\
A12 & Target market and pilot scenario require stakeholder validation. & Unresolved. & Interview supervisor/sponsor and at least one procurement/financing domain representative. \\
\bottomrule
\end{longtable}

\section{Business Rule Validation}

\begin{longtable}{P{0.34\linewidth}P{0.22\linewidth}P{0.32\linewidth}}
\toprule
Business rule & Status & Comment \\
\midrule
Protected workflows require authenticated server-derived actor context. & Validated. & Auth/session contract and server route registration use authenticated pre-handlers. \\
Eligibility controls must block protected procurement/escrow/PLS actions when parties are not eligible. & Validated for implemented flows; needs end-to-end pilot test. & Procurement, escrow, and PLS route registration include eligibility gateway behavior. \\
No raw KYC, commercial documents, delivery images, payment credentials, unrestricted contract text, or bank details may be written on-chain. & Validated by architecture; must be tested continuously. & Fabric architecture and blockchain contract preserve proof-only boundary. \\
Blockchain proof UI must not fabricate transaction IDs or verified states. & Validated by evidence; keep regression tests. & Release validation and PBI-457 evidence support this rule. \\
Escrow release cannot bypass accepted order, delivery evidence, eligibility, and dispute-free conditions. & Validated by contract; inspect tests during next hardening cycle. & Escrow contract defines release conditions and state restrictions. \\
Payment settlement cannot be claimed until adapter status confirms it. & Validated by design; production rail unresolved. & Current state is settlement instruction/sandbox/manual adapter only. \\
PLS cannot imply guaranteed profit or principal. & Validated by project claim boundary; requires formal Shariah/legal review. & Keep seedbed language conservative. \\
External integration calls require scoped credentials, signature/HMAC, idempotency, and audit logging. & Validated by PBI-458 evidence; needs security review. & Next phase should add threat-model and rate-limit hardening. \\
\bottomrule
\end{longtable}

\section{Requirement and System Analysis}

\subsection{Actors and user goals}

The active user-facing actors are Administrator, Compliance Reviewer, Buyer/Procurement Officer, SME/Supplier, Shariah Reviewer, Bank/Financier, Auditor, Regulator/Reporting User, Security Operator, Platform Operator, and external integration clients. Their core goals are governance, eligibility review, procurement execution, delivery proof, escrow/payment control, Shariah/PLS governance, audit/export verification, operational monitoring, and external event ingestion.

\subsection{Primary business process}

The main process is:

\begin{enumerate}
  \item User authenticates with credential-only login.
  \item Platform resolves trusted actor context and role/organization access.
  \item Compliance/eligibility gates supplier and buyer participation.
  \item Buyer and supplier negotiate or reference machine-readable terms.
  \item Buyer creates order; supplier acknowledges or rejects.
  \item Supplier or external device/client submits delivery evidence.
  \item Escrow evaluates release conditions, hold/dispute/arbitration, and settlement instruction readiness.
  \item Payment adapter records sandbox/manual instruction and reconciliation status.
  \item Lifecycle events are hashed and anchored when proof infrastructure is available.
  \item Auditor/regulator verifies proof/export bundle.
\end{enumerate}

\subsection{UML diagram sources}

The following PlantUML files are included with this report:

\begin{itemize}
  \item \repo{docs/uml/business-study-use-case.puml}
  \item \repo{docs/uml/business-study-procure-to-settlement-activity.puml}
  \item \repo{docs/uml/business-study-domain-class.puml}
  \item \repo{docs/uml/business-study-escrow-state-machine.puml}
\end{itemize}

They should be rendered during the documentation build or reviewed directly as version-controlled diagram source.

\section{External Research Findings}

\begin{longtable}{P{0.24\linewidth}P{0.46\linewidth}P{0.18\linewidth}}
\toprule
Source & Finding used in validation & Design implication \\
\midrule
Hyperledger Fabric docs & Fabric is an enterprise-grade permissioned DLT with known participants and modular architecture. & Fabric is appropriate for proof anchoring/consortium evidence but should not replace operational workflow persistence. \\
Hyperledger Fabric private data docs & Private data collections keep private data with authorized organizations and write hashes to the channel ledger. & Proof/hash-only and off-chain raw data rules are consistent with Fabric design. \\
GS1 EPCIS 2.0 & EPCIS defines visibility events and standard interfaces for capture/query/share of visibility event data. & IoT/QR/logistics support should stay as adapter/interface mapping rather than hard-coded hardware logic. \\
OCDS & OCDS release structures cover parties, planning, tender, award, contract, and implementation sections and treat releases as immutable events. & Contract/export mapping should use release-package semantics for disclosure/export, not replace the internal contract model. \\
ISO 20022 & The ISO 20022 repository includes a data dictionary, business-process catalogue, message definitions, and payments business areas. & Internal payment instructions can map to ISO-20022-like artifacts, but bank certification remains outside current evidence. \\
\bottomrule
\end{longtable}

\section{Deviation Analysis}

\begin{longtable}{P{0.34\linewidth}P{0.24\linewidth}P{0.30\linewidth}}
\toprule
Deviation / risk & Severity & Required action \\
\midrule
Newer modules still depend on in-memory adapters in runtime composition while deployable model uses Postgres for only part of the domain. & High for pilot. & Add Postgres migrations/adapters for documents, contracts, payments, PLS, certificates, integration credentials/idempotency, incidents, and export signing metadata. \\
PBI-462 remains planned; standards mapping is not closed. & High for product-market and interoperability. & Run formal standards mapping spike before production ERP/OCDS/UBL/Peppol claims. \\
PBI-438 remains planned; Fabric consortium is templates only. & High for blockchain production claims. & Build production-like Fabric lab with real MSP/CA/channel lifecycle evidence. \\
Some process docs still mention Aider/Sprint 1 assumptions. & Medium. & Refresh durable process docs to current Codex/agent-neutral and production-extension context. \\
Container model uses local development secrets/defaults and no managed TLS/secrets. & High for deployment. & Add environment hardening, secrets policy, image scanning, non-root containers, TLS/reverse proxy guidance. \\
Legal/Shariah/compliance stakeholder validation is not recorded. & High for external pilot. & Prepare review pack and log decisions before external pilot or investor-facing claims. \\
\bottomrule
\end{longtable}

\section{Risks and Open Questions}

\begin{enumerate}
  \item Which jurisdiction and retention policy is the first deployable pilot targeting?
  \item Which actors are legally authorized to approve Shariah certificate artifacts, export signing profiles, and dispute/arbitration outcomes?
  \item Should the next release prioritize full Postgres persistence or production-like Fabric lab execution first?
  \item Which external integration is most valuable for the first pilot: ERP/accounting, EPCIS/logistics, or payment sandbox?
  \item What is the acceptable minimum operational security baseline: local Docker pilot, VPS pilot, or managed cloud pilot?
  \item What data should be retained, redacted, encrypted, or legally held for KYC, documents, audit logs, and export bundles?
\end{enumerate}

\section{Recommendations}

\subsection{Immediate next phase}

Begin a validation-and-hardening phase rather than a new feature phase. The recommended next cycle is:

\begin{enumerate}
  \item Close PBI-462 as a standards mapping and business-rule validation spike.
  \item Produce a persistence coverage matrix and implement missing Postgres adapters/migrations for pilot-critical modules.
  \item Run a live deployable environment rehearsal using \repo{docker-compose.app.yml}, seeded accounts, browser smoke, backup/restore dry run, and evidence capture.
  \item Create a Fabric lab plan for PBI-438 but do not close PBI-438 until real CA/MSP/channel/lifecycle and cross-org smoke evidence exists.
  \item Update stale durable docs so process, architecture, and contract language match current production-extension state.
\end{enumerate}

\subsection{Frontend framework consideration}

Do not change frontend framework now. The existing React/Vite web application is sufficient for the current internal-pilot target. React Native is not justified until mobile-first requirements, offline capture, camera/QR scanning, field logistics workflows, and device-management constraints are validated. If QR scanning becomes a hard user requirement, evaluate responsive web/PWA first because it preserves the current codebase and deployment model.

\section{Next Development Plan}

\begin{longtable}{P{0.16\linewidth}P{0.42\linewidth}P{0.30\linewidth}}
\toprule
Milestone & Objective & Acceptance evidence \\
\midrule
M1 Standards and assumptions closure & Close PBI-462; map OCDS, UBL/Peppol, EPCIS, ISO 20022, Fabric, and Shariah rules to internal fields and boundaries. & Standards mapping report, traceability table, updated contracts, external citations. \\
M2 Persistence completion & Replace pilot-critical in-memory repositories with Postgres adapters/migrations. & Migration/seed tests, repository tests, data survival after restart. \\
M3 Deployable pilot rehearsal & Run containerized app, live DB, browser smoke, backup/restore dry run, and security/ops checks. & Evidence file with commands, environment, screenshots if applicable, blockers. \\
M4 Fabric lab validation & Build production-like Fabric lab and smoke-test AuditAnchor with real MSP/channel lifecycle. & Package IDs, approvals, commit, verified/mismatch/notFound smoke results. \\
M5 Security/compliance hardening & Add secrets policy, role matrix review, retention/legal hold plan, incident response and audit-event catalog. & Threat model, control checklist, updated runbooks, regression tests. \\
M6 Stakeholder validation & Validate business process with supervisor/sponsor/domain stakeholders. & Interview notes, accepted assumptions, revised backlog priorities. \\
\bottomrule
\end{longtable}

\section{Appendix A: Evidence and Citation Table}

\begin{longtable}{P{0.32\linewidth}P{0.50\linewidth}}
\toprule
Evidence & Use \\
\midrule
\repo{docs/evidence/qa/PRODUCTION_EXTENSION_RELEASE_VALIDATION.md} & Current release state, completed/open production-extension rows, browser smoke, and claim boundaries. \\
\repo{backlog/production-extension-roadmap.csv} & PBI-436 to PBI-462 status and dependencies. \\
\repo{src/app/server.ts} & Runtime module registration and partial Postgres runtime composition. \\
\repo{docker-compose.app.yml} & Internal deployable stack definition. \\
\repo{docs/architecture/PRODUCTION_FABRIC_CONSORTIUM_ARCHITECTURE.md} & Fabric consortium architecture and PBI-438 completion boundary. \\
\url{https://hyperledger-fabric.readthedocs.io/en/latest/whatis.html} & Fabric permissioned, modular enterprise DLT reference. \\
\url{https://hyperledger-fabric.readthedocs.io/en/latest/private-data/private-data.html} & Fabric private data collection and ledger hash behavior. \\
\url{https://standard.open-contracting.org/latest/en/schema/reference/} & OCDS release structure and contracting-process sections. \\
\url{https://ref.gs1.org/standards/epcis/2.0.0/} & GS1 EPCIS 2.0 visibility event and capture/query interface reference. \\
\url{https://www.iso20022.org/iso-20022-message-definitions} & ISO 20022 message-definition repository and payment business-area reference. \\
\bottomrule
\end{longtable}

\section{Appendix B: Raw Notes}

\begin{itemize}
  \item Current system is not assumed correct by default. This report treats implementation behavior as a hypothesis and records validation status per assumption.
  \item The largest architecture risk is not missing UI capability; it is persistence, operational hardening, and externally validated business-rule completeness.
  \item UML files are intentionally stored as PlantUML source so they can be reviewed and version controlled.
  \item This report does not run runtime commands. It records repository inspection and research findings. Independent validation should run build/test/deploy commands before merge into \repo{main}.
\end{itemize}

\end{document}
